% !TEX program = xelatex
\documentclass[11pt,a4paper]{article}

% ========== パッケージ群 ==========
\usepackage{amsmath}
\usepackage{amssymb}
\usepackage{amsfonts}
\usepackage{graphicx}
\usepackage{cite}
\usepackage{url}
\usepackage{xcolor}
\usepackage{fontspec}
\usepackage{xeCJK}

% ========== 日本語フォント設定 ==========
% XeLaTeX 前提。環境ごとに存在する日本語フォントが異なるため、
% TeX Live 標準の Harano Aji、Noto CJK、macOS 標準フォントの順にフォールバックする。
\defaultfontfeatures{Ligatures=TeX}
\IfFontExistsTF{HaranoAjiMincho-Regular}{%
  \setCJKmainfont{HaranoAjiMincho-Regular}[BoldFont=HaranoAjiGothic-Medium]
  \setCJKsansfont{HaranoAjiGothic-Regular}[BoldFont=HaranoAjiGothic-Bold]
}{%
  \IfFontExistsTF{Noto Serif CJK JP}{%
    \setCJKmainfont{Noto Serif CJK JP}[AutoFakeBold=2.0]
    \setCJKsansfont{Noto Sans CJK JP}[AutoFakeBold=2.0]
  }{%
    \IfFontExistsTF{Hiragino Mincho ProN}{%
      \setCJKmainfont{Hiragino Mincho ProN}[AutoFakeBold=2.0]
      \setCJKsansfont{Hiragino Sans}[AutoFakeBold=2.0]
    }{%
      \IfFontExistsTF{Yu Mincho}{%
        \setCJKmainfont{Yu Mincho}[AutoFakeBold=2.0]
        \setCJKsansfont{Yu Gothic}[AutoFakeBold=2.0]
      }{%
        \setCJKmainfont{IPAMincho}[AutoFakeBold=2.0]
        \setCJKsansfont{IPAGothic}[AutoFakeBold=2.0]
      }%
    }%
  }%
}
\xeCJKsetup{CJKmath=true}

\usepackage{hyperref}
\usepackage{geometry}
\geometry{margin=1in}

% ========== 定理環境 ==========
\usepackage{amsthm}

\theoremstyle{definition}
\newtheorem{lemma}{補題}[section]
\newtheorem{definition}{定義}[section]
\newtheorem{proposition}{命題}[section]
\newtheorem{assumption}{仮定}[section]
\newtheorem{remark}{注記}[section]
\newtheorem{phenomenological}{現象論的原理}[section]
\newtheorem{clarification}{説明}[section]

% ========== UTF-8 対応（pdfLaTeX） ==========
% macOS での使用時は XeLaTeX で実行してください
% 本ファイルはコンテナ環境での検証用（日本語フォント設定省略）

% ========== ハイパーリンク設定 ==========
\hypersetup{
  colorlinks=true,
  linkcolor=blue,
  citecolor=blue,
  urlcolor=blue,
  pdftitle={意味の非収束と共鳴条件},
  pdfauthor={加納 智之},
  pdfsubject={Resonanceverse 理論の基礎},
  pdfkeywords={共鳴, 誤配, 意味動力学, 非収束安定性},
}

% ========== バージョン管理 ==========
\newcommand{\versionnum}{1.2}
\newcommand{\versiondate}{2026/05/03}

% ========== タイトル情報 ==========
\title{意味の非収束と共鳴条件：\\
誤配に基づく Resonanceverse の基礎理論\\
\large \textit{Version \versionnum, \versiondate}}

\author{加納 智之$^{1,2,*}$}

\date{\today}

% ========== 本文開始 ==========
\begin{document}

\maketitle

% ========== 所属・連絡先（タイトル直下） ==========
\begin{center}
\small
$^{1}$ ZYX Corp 株式会社 人工叡智研究室（Artificial Sapience Laboratory）、名古屋、日本\\
$^{2}$ 独立研究者（Independent Researcher）\\
$^{*}$ 責任著者（Corresponding author）\\
E-mail: \url{tomyuk@zyxcorp.jp}\\
ORCID: \url{https://orcid.org/0009-0004-8213-4631}\\
\vspace{0.3cm}
\textbf{バージョン \versionnum} \quad \versiondate
\end{center}

\vspace{0.5cm}

\begin{abstract}
本論文は Resonanceverse を提示する。これは意味を誤配（misdelivery）に条件付けられた非収束的動力学プロセスとして再定義する理論的枠組みである。ポスト構造主義、特に Jacques Derrida は意味の不安定性を開示したが、そうした洞察は動力学モデルと結合されない限り、概ね記述的なままである。

従って、意味は $\mathbb{R}^d$ 内の符号化された状態軌道としてモデル化される。ここで差異と規範が十分に定義されている。形式的には、$\Delta M_{ij}(t) := M_j(t) - M_i(t)$ が二つのエージェント $(i, j)$ 間の意味論的状態差を時刻 $t$ で表す。ここで $M_k(t) \in \mathbb{R}^d$ である。

誤配は $\Delta M$ と同一ではない。むしろそれは、相互作用論的な生成的層として機能し、通じてそうした差異が生成され、持続し、チャネルを通じて再形成される。それでも $\Delta M$ は、状態がいったん符号化されたら、定量的層でどのように不整合が進化するかをインデックスする。

関係は文脈的重みと差異を結合する時間積分としてモデル化される。共鳴には意味的、不確実性、および価値制約 $C_M(t), C_U(t), C_V(t)$ の同時満足が必要である。本論文は、意味的アイデンティティへの収束なしに、実用的安定性がいかにして現れるかを明確にする。

本アプローチは意味を固定的な実体ではなく、継続的な偏差の持続的でかつ構造的に制約された過程として再定義する。ポスト構造主義哲学と動力学系理論を架橋することにより、Resonanceverse は人工知能、通信理論、制度設計への応用のための形式的基礎を提供する。

数値シミュレーションは、決定的な経験的証明ではなく構成的例示として用いられる。Phase 1～4 の実験は、有界な非収束レジーム、制度的制約の効果、共鳴に類する制約満足、および反証可能性を志向した診断が、明示的に指定された動力学の下で生成されうることを示す。これらのシミュレーションは本枠組みの内的整合性を支持するが、それ自体によって非収束公理を普遍的な経験法則として確立するものではない。
\end{abstract}

\section*{キーワード}
共鳴、誤配、意味動力学、非収束安定性、マルチエージェント $\Delta M$ ネットワーク、コンセンサス、共鳴条件、ポスト構造主義、ジャック・デリダ、動力学系、制度的制約、方法論的類推、Resonanceverse

\section{導入}

意味とは何か。この問いは哲学、言語学、認知科学、情報科学において何度も提起されてきた。しかし多くの理論は、明示的であれ暗黙的であれ、意味を比較的に安定した構造として、あるいは最終的に同定可能な状態として前提している。

この前提はポスト構造主義思想によって根本的に再検討された。とりわけ Jacques Derrida は、意味が自己同一性を持たず、\textit{différance} として延期され、散在することを示し、それゆえその不確定性と非閉鎖性を開示した。東浩紀はさらにこの不確定性を「郵便的脱構築」として再解釈し、意味が誤配を通じて生成されるというモデルを提案した。

しかし、こうした議論は概ね記述的・批判的枠組みに留まっている。それらは、意味の時間的変換と相互作用的動力学を記述できる形式的理論にまで至っていない。言い換えれば、以下の問いは十分に形式化されていない。

\textit{意味はいかにして時間の中で変化し、持続し、崩壊するのか？}

本研究は意味動力学の枠組みを導入することでこの隙間に対処する。意味を時間的に変動する状態量として扱う。形式的モデルでは、各エージェント $i$ に対して意味状態 $M_i(t)$ が割り当てられ、その差異
\[
\Delta M_{ij}(t) := M_j(t) - M_i(t)
\]
は二つのエージェント間の意味論的状態差として扱われる。この $\Delta M$ は誤配そのものと同一ではない。むしろ、チャネル全体での偏差、変位、再配送といった誤配プロセスが $\Delta M$ を時間的に生成し、増幅し、あるいは減衰させる。誤配プロセスは実用的な生成的層として機能し、一方 $\Delta M$ はそのプロセスが動力学系にマップされたときの基本的な観測量として機能する。このポジションはエラーを単なるノイズとして扱うアプローチとは対照的である。

\subsection{中心的仮説}

本論文の中心的仮説は以下の通りである。

\textit{意味は正しく伝達されるから持続するのではなく、継続的に誤配されるがゆえに非収束的な動力学的構造として持続する。}

この視点から、意味は静的なオブジェクトではなく、関係的かつ時間的なプロセスとして定義される。意味は単一のエージェント内部に存在しない。むしろそれは複数のエージェント間の相互作用を通じて生成される。従って、意味変化はエージェント間の差異のネットワーク、すなわち $\Delta M$ ネットワークとして記述される。

これは従来のコンセンサス形成理論に仮定されたコンセンサスの概念の再検討を要求する。本論文ではコンセンサスを意味の同一性ではなく、エージェント間の $\Delta M$ の分布が有界で安定した構造を形成する状態として再定義する。この意味において、コンセンサスは偏差の消滅ではなく、偏差の構造的修正である。

さらに本論文は意味的安定性を収束ではなく、非収束下での有界性として再定義する。この枠組みでは、意味は原理的に収束せず、継続的に変動する。しかし、そうした変動が特定の構造的条件を満たすとき、システム全体は安定した振る舞いを示す。このような状態は\textit{共鳴}と呼ばれる。

\subsection{論文の貢献}

このフレームワークに基づいて、本論文の貢献は五つである。

\begin{enumerate}
\item 意味の非固定的かつ非収束的性質を公理系として形式化する。
\item 誤配プロセスを生成的層として位置付け、$\Delta M$ ネットワークの動力学を対応する動力学系の観測可能層として導入する。
\item コンセンサス形成をマルチエージェント $\Delta M$ ネットワークの形成として再定義する。
\item 共鳴条件を非収束下での安定性の条件として提示する。
\item 制度設計およびエージェント・アーキテクチャへの応用可能性を示唆し、説明的な数値シミュレーションにより補助する。
\end{enumerate}

本研究はポスト構造主義の洞察を否定しない。むしろ、そうした洞察を保存しながら、意味理論を記述から動力学系理論へ移動させることを目指している。この再定義は意味を固定的なオブジェクトではなく、関係の内部に持続する動的現象として扱う。それゆえ、インテリジェンスのモデル化、人工知能システムの設計、および社会制度の理論化のための新たな基礎を提供している。

\section{意味動力学の公理系}

本セクションは意味を時間的に変動する状態量として扱うための公理系を定義する。本研究では、意味は固定的な表現や内在的属性ではなく、エージェント間の相互作用を通じて継続的に更新される動的状態である。従って意味は単一のポイントへの収束として記述されるのではなく、時間を通じて移動する状態として記述される。

時刻 $t$ にエージェント $i$ が保有する符号化された意味状態を $M_i(t) \in \mathbb{R}^d$ で表す。$\mathcal{M}$ を抽象的な意味状態空間とする。測定可能な状態は符号化マップ
\[
\phi : \mathcal{M} \to \mathbb{R}^d
\]
を通じて $\mathbb{R}^d$ で表現されていると仮定される。

記述の簡潔性のため、符号化された像を単に $M_i(t)$ で表記する。差異と規範は $\mathbb{R}^d$ の通常の線形代数的構造を用いて定義される。距離は有するが減法をもたない一般的な計量空間では $\Delta M_{ij} = M_j - M_i$ を記述しない。

二つのエージェント $i$ と $j$ 間の符号化された状態差は次のように定義される。
\[
\Delta M_{ij}(t) := M_j(t) - M_i(t) \in \mathbb{R}^d.
\]

形式的には、$\Delta M_{ij}(t)$ は二つのエージェントが保有する意味状態間の差異である。従って $\Delta M$ は誤配そのものと同定されるべきではない。むしろ誤配はチャネル、制度、規範の実用的側にあることもある。本論文では、$\Delta M$ を誤配プロセスを通じて現れ、持続し、変換される意味的差異として制約付けで定義する。この区別は理論の数学的および修辞的構造の枢軸となる。

本論文では $\Delta M$ を多次元ベクトルとして扱う。その大きさは従ってノルムで測定される。明示されない限り、ユークリッド規範が使用される。
\[
\|\Delta M_{ij}(t)\| := \left(\sum_{k=1}^{d} (\Delta M^{(k)}_{ij}(t))^2\right)^{1/2}.
\]

各成分は意味、価値、不確実性など異なる次元に対応することがあり得る。

従来の理論では、そのような量は誤解またはノイズとして扱われることがある。本研究では、それを意味変化の基本的単位として扱う。意味は恒等性ではなく差異によって駆動される。そのような差異は例外的ではない。それは通常の条件であり、意味生成の出発点である。

\subsection{核心的公理}

この前提に基づいて、最初の公理は意味の時間的性質に関わる。全てのエージェントについて、意味状態は時間上で不変ではない。
\[
\forall I \subset [0,\infty), \mu(I) > 0 \Rightarrow \exists t_1, t_2 \in I \text{ such that } M_i(t_1) \neq M_i(t_2).
\]

この公理は意味が完全に静的なままではないことを意味する。正の測度を持つ任意の区間においても、意味は少なくともわずかに変動する。完全な恒等状態は時間全体を通じて持続しない。

さらに、意味生成の過程は孤立的ではなく、他者との相互作用に依存する。これは次のように表現される。
\[
M_i(t) = M_i(0) + \int_0^t \sum_j f(\alpha_{ji}(\tau), \Delta M_{ji}(\tau)) d\tau,
\]
ここで関数 $f$ は文脈的親和性と意味的差異をエージェント $i$ の意味状態の変化ベクトルへマップする。
\[
f : [0, 1] \times \mathbb{R}^d \to \mathbb{R}^d.
\]

この表式はエージェント $i$ の意味状態がその初期状態に他のエージェントとの関係で生じる意味的差異の累積を加えたものから構成されることを述べている。それは意味が内部生成ではなく関係的プロセスとして形成されることを明確にする。言い換えれば、意味は他者からの偏差の歴史である。


\subsection{非収束性公理}

次に、次の公理は意味的差異の時間的極限に関わる。
\[
\limsup_{t \to \infty} \|\Delta M_{ij}(t)\| > 0.
\]

この非収束性公理は、エージェント間の意味的差異が時間の中で消滅しないことを述べる。どれほど相互作用が継続しても、完全な恒等性は原理的に達成されない。この点により、本理論は意味の最終的統一を前提するモデルから区別される。

\begin{proposition}[スケール分離と両立性]
非収束性公理
\[
\limsup_{t \to \infty} \|\Delta M_{ij}(t)\| > 0
\]
は、軌道がゼロの小近傍へ入ることを禁止しない。それが否定するのは、差異の完全な消滅、すなわち厳密な意味的恒等性である。

より正確には、実用的スケール $\bar{\varepsilon}>0$ と時刻 $T\geq 0$ が存在し、
\[
0 < \|\Delta M_{ij}(t)\| < \bar{\varepsilon} \quad \text{for all } t \geq T,
\]
が成り立つ一方で、同時に
\[
\limsup_{t \to \infty} \|\Delta M_{ij}(t)\| > 0
\]
を満たすことは整合的である。この二層構造、すなわち論理的層における「ゼロへ向かわないこと」と、操作的層における「実用的スケールで修正可能であること」は、\textit{スケール分離}と呼ばれる。

非収束性公理は導出されるものではなく、置かれるものである。以下の仮定は、リアプノフ証明を過剰に主張することなく、非収束性と有界性の両立を保守的に表現する方法を与える。
\end{proposition}

\begin{assumption}[究極有界な意味動力学]
意味動力学は有界な吸収集合 $B \subset \mathbb{R}^d$ を持つと仮定する。すなわち、任意の許容初期条件に対し、ある $T \geq 0$ が存在して、
\[
M_i(t) \in B \quad \text{for all } t \geq T
\]
が成り立つ。
同値的に、次の散逸条件を仮定してもよい。proper なリアプノフ関数
$V:\mathbb{R}^d\to\mathbb{R}_{\geq 0}$ と定数 $R,\gamma>0$ が存在し、許容軌道に沿って、
\[
\|M_i(t)\| \geq R \quad \Rightarrow \quad \frac{d}{dt}V(M_i(t)) \leq -\gamma
\]
が成り立つ。この点ごとの散逸形式は、非負のリアプノフ関数の長期平均微分が十分大きな全ての時間幅で厳密に負であり続ける、というより強く一般には不整合になりうる主張を避ける。
\end{assumption}

\begin{proposition}[有界性と非収束性の両立性]
意味動力学が有界な吸収集合 $B \subset \mathbb{R}^d$ を持ち、関連する相互作用対 $(i,j)$ について非収束性公理
\[
\limsup_{t\to\infty}\|\Delta M_{ij}(t)\| > 0
\]
が成り立つと仮定する。このとき、有界性と非収束性は相互に両立する。すなわち、ある $T \geq 0$ と $\varepsilon < \infty$ が存在して、
\[
\|\Delta M_{ij}(t)\| < \varepsilon
\quad \text{for all } t \geq T,
\]
が成り立つ一方で、意味的恒等性は漸近的には到達されない。
\end{proposition}

\begin{proof}
吸収集合の仮定により、有界集合 $B$ と時刻 $T$ が存在し、全ての $t\geq T$ に対して $M_i(t),M_j(t)\in B$ である。したがって、
\[
\|\Delta M_{ij}(t)\| = \|M_j(t)-M_i(t)\| \leq \operatorname{diam}(B) < \infty
\quad (t\geq T)
\]
が成り立つ。$\varepsilon:=\operatorname{diam}(B)$ とおく。非収束性公理は、独立した条件
$\limsup_{t\to\infty}\|\Delta M_{ij}(t)\|>0$ を与える。この二つの主張は同時に満たされうる。前者は振幅を上から制約し、後者は厳密な恒等性への漸近的崩壊を防ぐ。これにより示されるのは両立性であり、公理の導出ではない。
\end{proof}

このような意味的差異の累積が、エージェント間の関係を構成する。関係は次のように定義される。
\[
R_{ij}(t) = \int_0^t g(\alpha_{ij}(\tau), \Delta M_{ij}(\tau)) d\tau,
\]
ここで $g : [0, 1] \times \mathbb{R}^d \to \mathbb{R}^r$ である。

関数 $g$ は、文脈的親和性と意味的差異を関係状態へ写像する。その値域は一般に $\mathbb{R}^r$ である。この定義の下では、関係は静的属性ではなく、過去にどのような意味的差異が生じ、それらがどのように累積されたかの履歴である。仮に二つのエージェントが現在同一の意味状態を共有していても、それぞれの関係履歴が異なれば区別されうる。

しかし、全ての意味的差異が同じように作用するわけではない。エージェント間の文脈的親和性は、意味変化の影響を調整する。この文脈的親和性を $\alpha_{ij}(t)$ と書き、$\alpha_{ij}(t) \in [0, 1]$ とする。

有効な意味変化は次のように定義される。
\[
\Delta M^{\text{eff}}_{ij}(t) = \alpha_{ij}(t) \Delta M_{ij}(t).
\]

この式は、意味的差異が意味更新へ直接寄与するのではなく、関係と文脈に応じて重み付けされることを示す。したがって意味変化は単なる差異の関数ではなく、関係に媒介された過程である。

意味状態の時間発展は次のように記述される。
\[
\frac{dM_i(t)}{dt} = \sum_j f(\alpha_{ji}(t), \Delta M_{ji}(t)).
\]

ここで $f$ は一般に非線形であり、履歴依存性を含むこともありうる。この定式化は、意味が単純な加算過程ではなく、複雑な相互作用システムとして振る舞うことを示す。

マルチエージェント・システム全体のレベルでは、この相互作用はネットワークとして表される。$G(t) = (V, E(t))$ を、エージェント集合 $V$ と相互作用集合 $E(t)$ からなるグラフとする。各エッジはエージェント間の意味的差異 $\Delta M_{ij}(t)$ に対応する。意味は個々のエージェントの内部に閉じ込められているのではなく、ネットワーク全体に分散している。

このようなシステムで重要なのは、意味的差異が消滅しなくても安定しうることである。本論文では安定性を次の条件として定義する。
\[
\exists T \geq 0, \exists \varepsilon > 0 \forall t \geq T, \forall (i,j) \in E(t) : 0 < \|\Delta M_{ij}(t)\| < \varepsilon.
\]

この条件は、意味的差異が存在し続けながらも、無限に発散するのではなく、一定範囲内に有界に保たれることを意味する。意味は収束せずに安定する。この非収束安定性は、収束に基づく古典的安定性概念とは異なる概念的枠組みを与える。

最後に、時間の中で消滅しない残余的意味差を導入する。グラフ $G(t)$ との混同を避けるため、残余は $G_{ij}(T)$ と書き、
\[
G_{ij}(T) := \sup_{t \geq T} \|\Delta M_{ij}(t)\|
\]
と定義する。

この量は、相互作用によって完全には解消されず、システムの長期的振る舞いに影響し続ける差異を表す。残余は、意味が決して完全には統一されないという事実の数学的表現であると同時に、過去の相互作用が現在へ影響し続ける機構を与える。


\subsection{説明的二エージェント例}

曖昧な符号規約に依存せずに公理的枠組みを説明するため、$\mathbb{R}^2$ における二エージェントシステムを考え、相対状態を直接
\[
x(t) := M_2(t)-M_1(t)
\]
と定義する。相対動力学は
\[
\dot{x}(t)=F(x(t))
\]
により指定される。ここで、
\[
F(x)=
\begin{cases}
\lambda_{\mathrm{rep}}x, & \|x\|<r,\\
-\lambda_{\mathrm{att}}x, & \|x\|\geq r,
\end{cases}
\]
であり、$\lambda_{\mathrm{rep}},\lambda_{\mathrm{att}}>0$ かつ $r>0$ とする。この相対的定式化では解釈が曖昧でない。半径 $r$ の内側では、エージェントは厳密な恒等性から押し離される。半径 $r$ の外側では、その差異は許容領域へ引き戻される。

滑らかな実装では、区分的な場を正則化された半径方向係数で置き換えることができる。たとえば、
\[
F(x)=\alpha(t)\left(\frac{r^2}{\|x\|^2+\epsilon}-1\right)x,
\]
であり、$\epsilon>0$ は原点での特異性を防ぐ。これは滑らかなリミットサイクルの存在を証明するものではなく、環状捕捉領域へ向かう傾向を生成する。したがって、この例はリミットサイクル存在に関する一般定理ではなく、有界な非収束性の構成的例示として読むべきである。

可能な数値設定の一つは次の通りである。
\begin{itemize}
\item 初期状態：$M_1(0) = [0.0, 0.0]$、$M_2(0) = [1.0, 0.5]$。
\item 初期差異：$x(0)=\Delta M_{12}(0)=[1.0,0.5]$、$\|x(0)\|\approx 1.12$。
\item 文脈的親和性：$\alpha(t)=0.7+0.2\sin(0.1t)$。
\item 目標半径：$r=0.1$。実用的な許容帯として、たとえば $[0.05,0.15]$ を置く。
\end{itemize}

期待される振る舞いは、大きな差異が縮小され、ゼロ近傍への崩壊が反発され、軌道が小さい非ゼロ帯の中に保たれることである。これはスケール分離を説明する。意味的差異は実用的には修正可能になりうるが、厳密な恒等性には到達しない。

この公理系を通じて、意味は個体の属性ではなく、エージェント間の差異とその履歴から構成される動的構造として形式化される。本セクションで導入された枠組みは、次のセクションにおいてマルチエージェント・システムにおけるコンセンサス形成を再定義する基礎となる。

\section{マルチエージェント $\Delta M$ ネットワークとコンセンサス形成}

本セクションはセクション2で定義された意味動力学をマルチエージェントシステムに拡張し、コンセンサス形成を再定義する。従来のコンセンサス理論はしばしば暗黙的に、エージェント間の意味が一致する状態、すなわち $M_i = M_j$ を仮定している。しかし、セクション2で示されたように、意味は非収束的であり、意味的差異 $\Delta M_{ij}(t)$ は原理的に消滅しない。

この前提の下では、完全な恒等としてのコンセンサスは一般的に成立しない。

この問題に対処するため、本研究はコンセンサスを意味の恒等性ではなく、意味的差異の構造として再定義する。コンセンサスは、指定された条件下で $\Delta M$ ネットワーク間のエージェントが安定した分布を形成する状態である。

エージェント集合を $\mathcal{A} = \{a_1, a_2, \ldots, a_n\}$ とする。すなわち $\mathcal{A}$ は意味的相互作用ネットワークに参加するエージェントの有限集合を表す。

各エージェントは符号化された状態 $M_i(t) \in \mathbb{R}^d$ を持つ。意味的差異がその成分である構造行列 $X(t)$ を $X(t)$ で表す。マトリックスレンダリングの不安定性を避けるため、成分ごとに定義される。
\[
[X(t)]_{ij} := \Delta M_{ij}(t) \quad (i \neq j), \quad [X(t)]_{ii} := 0.
\]

この構造 $X(t)$ はマルチエージェントシステムの意味的差異ネットワークを表現する。ここで重要な点は意味が個々のエージェント内に限定されず、この差異のネットワークとして存在することである。意味は分散し、エージェント間の関係を通じてのみ定義される。

ネットワークの時間進化はセクション2で導入された動力学に従う。
\[
\frac{dM_i(t)}{dt} = \sum_j f(\alpha_{ji}(t), \Delta M_{ji}(t)).
\]

このインタラクションを通じて、意味的差異の分布は時間の中で変化するが、必ずしもゼロに収束しない。実際のマルチエージェントシステムでは、持続的な意味的差異は一般に規則ではなく例外である。

この視点に基づいて、コンセンサスは以下のように定義される。

\[
\text{コンセンサス}(t) \Leftrightarrow \forall (i,j) \in E(t) : \|\Delta M_{ij}(t)\| < \varepsilon \land -\delta < \frac{d}{dt}\|\Delta M_{ij}(t)\| < \delta.
\]

許容値 $\varepsilon$ はセクション2の非収束性公理 $\limsup_{t \to \infty} \|\Delta M_{ij}(t)\| > 0$ と矛盾しない。むしろそれは論理的層で差異の完全な消滅を除外しつつ、実用的許容帯を標示する。これはスケール分離として前に議論されたものである。

ここで $\varepsilon > 0$ と $\delta > 0$ は許容閾値である。後者の制約は一方向的条件 $\frac{d}{dt}\|\cdots\| < \delta$ ではなく、上記に示された二方向不等式である。それは規範の時間導関数が $-\delta$ と $\delta$ の間に留まることを要求し、それゆえ急速な増加と急速な減少の両方を抑制する。

この構造は分散の観点からも記述されることがある。平均エッジ規範を
\[
\mu_\Delta(t) = \frac{1}{|E(t)|} \sum_{(i,j) \in E(t)} \|\Delta M_{ij}(t)\|
\]
とする。

エッジ上の規範の分散は
\[
\text{Var}_E(\Delta M(t)) = \frac{1}{|E(t)|} \sum_{(i,j) \in E(t)} (\|\Delta M_{ij}(t)\| - \mu_\Delta(t))^2
\]
である。

この量はマルチエージェントシステムにおける意味的差異規範の広がりを表現する。分散が大きい場合、意味的差異は広く分散し、システムは断片化されていると解釈されることがある。分散が小さすぎる場合、意味的差異はほぼ不在であり、システムは均質化される。そのような状態は安定して見えるかもしれないが、外部変化に対して脆弱である。

従ってコンセンサスはゼロへ向かう分散ではなく、安定した中間範囲によって特徴付けられる。この状態は意味的差異が存在し続ける一方で、それらの変動が制御されることを意味する。

このように定義されたコンセンサスは従来の恒等性モデルと異なり、差異の存在を前提とする。コンセンサスは偏差の消滅ではなく、偏差の構造的修正である。

最後に、この枠組みにおける制度の役割に対処する必要がある。制度はエージェント間の相互作用を制約し、それゆえ $\Delta M$ の分布を制御する。制度は意味そのものを固定しない。それはエージェント間の差異の範囲と変動を制約する。

\begin{definition}[制度]
マルチエージェント $\Delta M$ ネットワークに作用する\textit{制度} $\mathcal{I}$ は三つ組 $\mathcal{I} = (\mathcal{G}, \mathcal{C}, \Pi)$ である。ここで：
\begin{itemize}
\item $\mathcal{G} \subseteq V \times V$ は規制対象のエージェント対の集合（制度的制御の範囲）；
\item $\mathcal{C} = \{c_{ij} : \mathbb{R}^d \to \{0, 1\}\}_{(i,j) \in \mathcal{G}}$ は制約述語の集合であり、$c_{ij}(\Delta M_{ij}) = 1$ は遵守を示す；
\item $\Pi : \mathbb{R}^d \to \mathbb{R}^d$ は制約違反時に遵守を強制する射影演算子である。
\end{itemize}

制度的に規制された動力学は以下である：
\[
\frac{dM_i}{dt} = \Pi\!\left(\sum_j f(\alpha_{ji}(t), \Delta M_{ji}(t))\right),
\]
ここで $\Pi$ は全ての制約が満たされるとき恒等写像として作用し、そうでない場合は更新を実行可能集合に射影する。単純なインスタンス化は $c_{ij}(\Delta M) = \mathbf{1}[\|\Delta M\| < \varepsilon]$ であり、セクション3のコンセンサス閾値を生成する。
\end{definition}

この定義は制度を静的規範ではなく、ベクトル場そのものに作用する動的制御メカニズムとして扱う。制度設計はそこで $(\mathcal{G}, \mathcal{C}, \Pi)$ を選択することにより、結果としての $\Delta M$ 分布が制御され非収束安定構造が維持される条件を設計する問題となる。

従ってマルチエージェントシステムにおける意味はエージェント間の差異のネットワークとして記述され、コンセンサスはそのネットワークにおける安定した構造として現れる。本セクションで提示された枠組みは次のセクションで導入される共鳴条件を通じて非収束性下での安定性の、より精密な定義へ発展する。

\section{共鳴と非収束安定性}

本セクションはセクション2および3で定義された意味動力学とマルチエージェント $\Delta M$ ネットワークに基づいて、非収束下での安定性の条件として共鳴を導入する。従来の理論では、安定性はしばしば収束を通じて定義される。すなわち、意味状態が固定点に到達することによってである。本研究の非収束性前提の下では、そのような定義を適用することができない。従ってそれは意味が収束しないにもかかわらず安定した振る舞いを示す状態を定義することが必要である。

意味的差異 $\Delta M_{ij}(t)$ が時間上で消滅しない状態を考え、一方その大きさは一定範囲内で留まる。
\[
\Delta M_{ij}(t) \neq 0 \land \|\Delta M_{ij}(t)\| < \varepsilon.
\]

この条件は意味的差異が存在し続ける一方で無限に発散する代わりに制約されている状態を表現する。そのような状態では、エージェント間の意味は完全に一致することはないが、差異の範囲が安定しているため、システム全体は秩序立った振る舞いを示す。この安定性は収束ではなく有界性によって特徴付けられる。

しかし、この条件だけでは不十分である。意味的差異の単なる有界性は周期的振動または混沌的振る舞いを含むことがある。従ってその規範の時間導関数も実用的閾値によって制約される必要がある。
\[
-\delta < \frac{d}{dt}\|\Delta M_{ij}(t)\| < \delta.
\]

これはセクション3のコンセンサス定義に同じ $\delta$ である。それは規範の変化が実用的に受け入れ可能なほど遅く留まる状態を指定する。

これらの条件を満たす状態は\textit{非収束安定性}と呼ばれる。それは意味が固定点に収束することなく有界領域内で変動し続ける状態を指す。重要な点は、変動そのものが不安定を暗示しないこと。むしろ、安定性は変動が制約された範囲内で維持されるがゆえに、ちょうどそこで生じることがあり得ることである。

この非収束安定性をより精密に特徴付けるため、本論文は\textit{共鳴}の概念を導入する。記号 $C_M(t), C_U(t), C_V(t)$ は各時刻で真理値を割り当てる述語である。それらは意味状態ベクトル $M_i(t)$ の省略形ではない。以下にその形式的定義を与える。

\begin{definition}[制約関数 $C_M, C_U, C_V$]
$\theta_M, \theta_U, \theta_V > 0$ をドメイン固有の閾値とする。

\textbf{(i) 意味的分散制約。}
\[
C_M(t) :\Leftrightarrow \text{Var}_E(\Delta M(t)) < \theta_M,
\]
ここで $\text{Var}_E$ はセクション3で定義されたエッジ上の分散である。

\textbf{(ii) 不確実性整合制約。}
$U_i(t) \in [0,1]$ をエージェント $i$ の時刻 $t$ における正規化された認識論的不確実性とする。具体的には $U_i(t) = H(\pi_i(t)) / \log K$ であり、$H$ はシャノンエントロピー、$\pi_i(t)$ はエージェント $i$ の $K$ 状態上の信念分布、分母は $[0,1]$ への正規化である。以下を定義する：
\[
D_U(t) := \max_{(i,j) \in E(t)} |U_i(t) - U_j(t)|, \quad C_U(t) :\Leftrightarrow D_U(t) < \theta_U.
\]

\textbf{(iii) 価値方向制約。}
$V_i(t) \in \mathbb{R}^k$ をエージェント $i$ の評価ベクトル（例：効用関数勾配、価値優先順位）とし、$\|V_i(t)\| > 0$ とする。二つのエージェント間のコサイン距離を
\[
d_{\cos}(V_i, V_j) := 1 - \frac{V_i \cdot V_j}{\|V_i\|\,\|V_j\|} \in [0, 2]
\]
と定義する。価値乖離と制約は：
\[
D_V(t) := \max_{(i,j) \in E(t)} d_{\cos}(V_i(t), V_j(t)), \quad C_V(t) :\Leftrightarrow D_V(t) < \theta_V.
\]
\end{definition}

共鳴の概念そのものは論理積である。
\[
\text{共鳴}(t) \Leftrightarrow C_M(t) \land C_U(t) \land C_V(t).
\]

三つの制約は論理的に独立である：各々が異なる次元（意味的広がり、認識論的整合、評価的一貫性）に対処する。それらの同時満足は構造的情報の損失なしに単一のスカラ条件に還元できない。

\begin{remark}[零評価ベクトルを持つエージェント]
コサイン距離 $d_{\cos}$ は $\|V_i\| = 0$ のとき未定義である。そのようなエージェント——評価方向を持たないエージェント——は $C_V$ の評価から除外される。形式的には、$D_V(t)$ は制限されたエッジ集合 $E_V(t) := \{(i,j) \in E(t) : \|V_i(t)\| > 0 \text{ かつ } \|V_j(t)\| > 0\}$ 上で計算される。$E_V(t) = \emptyset$ の場合、条件 $C_V(t)$ は空虚に真である。
\end{remark}

\begin{definition}[共鳴]
\textit{マルチエージェントシステムが時刻 $t$ で共鳴状態にあるとは、意味論的分散条件（$C_M$）、不確実性整合条件（$C_U$）、および価値方向条件（$C_V$）が同時に満たされる場合である。共鳴は差異が複数の次元に渡って首尾一貫して制約される状態である。}
\end{definition}

ここで重要なのは、整合がここで完全な恒等を意味しないことである。$C_M$ は意味的差異が消滅されるのではなく、分散と閾値の観点から制御されることを要求する。同様に、$C_U$ は共有された不確実性の度合いを指し、$C_V$ は評価的標準または判断方向の整合を指す。

従って共鳴は差異が持続する一方で複数次元に渡って首尾一貫して制約される状態である。そのような状態では、意味的差異は混沌的に拡散しないが、特定の構造に従って分布する。この構造は意味状態空間の領域として表現されることがある。
\[
M_i(t) \in B \subseteq \mathbb{R}^d.
\]

ここで $B$ は符号化空間 $\mathbb{R}^d$ の有界領域であり、共鳴状態にあるエージェントの符号化された意味値はこの領域内に分布すると見なされる。この領域は固定点ではない。それは意味が動力学的に変動する範囲を指定する。

動力学系の視点から、そのような領域はアトラクタに類似している。しかし本研究のアトラクタは収束の点ではなく、非収束軌道が制約される領域である。この意味において、共鳴は\textit{非収束アトラクタ}として解釈されることがある。

意味的差異の長期的振る舞いを考慮して、消滅しない長期振幅が残る。規範を最初に取り、次を定義する。
\[
G^{(\infty)}_{ij} := \limsup_{t \to \infty} \|\Delta M_{ij}(t)\|.
\]

非負スカラ写像 $h_{ij}(t) = \|\Delta M_{ij}(t)\|$ に対して、上限の一般的公式は以下を与える。
\[
G^{(\infty)}_{ij} = \lim_{T \to \infty} G_{ij}(T) = \inf_{T \geq 0} \sup_{t \geq T} \|\Delta M_{ij}(t)\|.
\]

有限時刻 $t$ 付きの $\limsup$ を混合する早期ドラフト問題を避けるため、振幅は常に $\|\cdot\|$ と $h_{ij}$ を通じて記述される。

この残存は過去の相互作用によって生成された差異が完全には消滅せず、システム全体に影響し続けることを示す。共鳴状態では、この残存も一定範囲内で制約されている。共鳴は従って現在の状態だけではなく、その歴史を含む全体的構造の安定性を指す。

従って共鳴は非収束下での安定性の条件として位置付けられる。意味は収束することなく構造的に安定のままであることがある。この安定性は制約 $C_M, C_U, C_V$ の同時作用を通じて生じる。

共鳴はセクション3で導入されたコンセンサスの概念を一般化する。コンセンサスは主として意味的差異の大きさと安定性に焦点を当てるのに対し、共鳴は不確実性と価値の次元を追加的に含む。従ってコンセンサスは共鳴の特殊ケースとして理解されることができる。

\subsection{非収束安定性の理論的性質}

本小セクションは、セクション2～4が依拠する数学的基礎を簡潔に列挙する。完全な証明はここでは発展させない。厳密な発展は将来の研究に残される。「定理」という用語は避けられ、さらに仮定を要する主張は明示的に命題または現象論的原理として標識される。特に、更新項 $f$ の有界性だけではそれ自体は軌道の有界性を暗示しない。本小セクションの中心的仮定は従って正不変領域の存在である。

\begin{remark}[有界 $f$ だけは有界軌道を暗示しない]
更新関数 $f$ が一様有界であっても、
\[
M_i(t) = M_i(0) + \int_0^t \sum_j f(\alpha_{ji}(\tau), \Delta M_{ji}(\tau)) d\tau
\]
は一般的には有界である必要がない。$j$ 上の和が長時間にわたって時間依存的な累積で作用する場合、線形成長が生じることがある。従ってそれは有界振幅の安定領域を主張する場合、散逸、射影、または正不変性といったベクトル場についての追加的仮定が不可欠である。

$f$ の一様有界性、
\[
\|f(\alpha, x)\| \leq C \quad \text{for all } \alpha \in [0,1], x \in \mathbb{R}^d,
\]
はただ局所更新が任意に大きくならないことを意味するために取られている。
\end{remark}

\begin{assumption}[正不変有界閉領域]
時間的ベクトル場に対して有界閉集合 $B \subseteq \mathbb{R}^d$ が存在して、容認可能な初期条件から開始する全ての軌道が、一度 $B$ に入れば、永遠に $B$ に留まる。

幾何学的には、境界 $\partial B$ に内向的指示条件を課すことができ、例えば
\[
\left\langle \sum_j f(\alpha_{ji}, \Delta M_{ji}), n_i \right\rangle \leq 0
\]
外向き法線 $n_i$ に対して。本論文は上記の集合水準仮定を採用する。
\end{assumption}

\begin{proposition}[非収束だがかつ振幅有界なレジームの存在]
仮定1が成立し、非収束性公理
\[
\limsup_{t \to \infty} \|\Delta M_{ij}(t)\| > 0
\]
が効力を持つ場合、意味状態は有界領域内に制約されながら厳密な恒等に到達しないことがある。すなわち、有界領域 $B$ と時刻 $T$ が存在して
\[
M_i(t) \in B \quad (\forall t \geq T), \quad \limsup_{t \to \infty} \|\Delta M_{ij}(t)\| > 0.
\]

本命題は意図的に控えめである。これは、明示的な仮定の下で有界性と非収束性が両立しうることを述べるに留まり、非収束性公理を一般的な滑らかな動力学から導出するものではない。
\end{proposition}

\textbf{十分なモデル化条件。}
具体的モデルは、次の要素によって上の命題を実現しうる。
\begin{enumerate}
\item[(SC1)] \textbf{局所正則性：} 更新場が容認可能領域上で局所リプシッツ連続であり、軌道が適切に定義されること。
\item[(SC2)] \textbf{有界な容認可能領域：} 正不変な有界集合 $B$ が存在すること。これは例えば射影、飽和、または境界上の内向き条件によって実装されうる。
\item[(SC3)] \textbf{反崩壊機構：} 考察対象となる相互作用ペアについて、ベクトル場が厳密な恒等集合 $\Delta M_{ij}=0$ から反発する、周回する、またはその他の仕方でそれを避けること。
\end{enumerate}
これらのモデル化条件の下で、軌道は有界に留まりながら、厳密な意味的同一性への崩壊を避けることができる。

\textbf{例：相対的二エージェント実装。}
$x(t)=M_2(t)-M_1(t)$ とし、相対動力学を
\[
\dot{x}=F(x)
\]
として指定する。ここで
\[
F(x)=
\begin{cases}
\lambda_{\mathrm{rep}}x, & \|x\|<r,\\
-\lambda_{\mathrm{att}}x, & \|x\|\geq r,
\end{cases}
\]
であり、$\lambda_{\mathrm{rep}},\lambda_{\mathrm{att}}>0$ とする。この区分的例は、滑らかなリミットサイクルの存在を証明するものではなく、環状捕捉領域を生成するものとして解釈されるべきである。微分可能な例が必要な場合は、不連続な切替を滑らかな半径方向係数で置き換えればよい。

\textbf{証明方針。}
有界性は、$B$ への正不変性または散逸的射影から従う。非収束性は、反崩壊機構または非収束性公理そのものによって与えられる。その合成により、$\omega$ 極限挙動が厳密な恒等の単一点ではない振幅有界な軌道が得られる。任意のマルチエージェントネットワークに対する一般的なスペクトル判定基準は将来の研究に残される。

\textbf{解釈。} 意味が単一のポイントに収束しなくても、正不変領域は $\Delta M$ が内部で非零に留まりながら振幅を制約することができる。二つの層は共存できる。

\textbf{スケッチ。} 正不変性は $\|M_i\|$ の爆発を除外する。一方、非収束性公理は $\|\Delta M\| \to 0$ である軌道を除外する。典型的ケースは単調に収束しない閉または引力的運動を含む。

\begin{assumption}[有限次数]
相互作用グラフ $G(t)$ は有限次数を持つ。すなわち、与えられたエージェントに同時に影響するエージェント数は上から有界である。
\[
\deg(i,t) \leq D \quad \text{for all } i, t.
\]
\end{assumption}

\begin{phenomenological}[疑似収束からの乖離]
非収束システムでは、全ての単一のポイント $M^*$ が正真正銘の勾配収束の極限として機能するわけではない。小さな差異と擾乱は時間とともに累積され、状態は候補ポイント付近に留まるかもしれないが、最終的にそれから乖離する。これは証明された定理ではなく、行動的テンプレートとして提示されている。

\textbf{解釈。} 意味が静的に見えても、公理 $\limsup \|\Delta M\| > 0$ の下では、それは完全な恒等に沈むことはできず、運動に戻ることがある。

\textbf{スケッチ。} 非零差異からの漂流または内部動力学と外部制約を比較する。
\end{phenomenological}

\begin{assumption}[擾乱後の拡張領域]
外部決定論的ノイズ $\eta(t)$ または有界増分を持つ確率過程が状態に単調に作用する擾乱されたシステムに対して、システムは有界領域 $B' \subseteq \mathbb{R}^d$ で正不変に留まると仮定する。ここで $B'$ は元々の $B$ より大きい。これは入力と状態制約に問題を削減する。
\end{assumption}

\begin{phenomenological}[外部擾乱の吸収]
仮定3の下で、擾乱は異常として現れる可能性は少なく、むしろ既存の変動の上に層状化される。軌道はグローバルに $B'$ で制約され、共鳴としての構造は閾値超過によってのみ評価される。
\end{phenomenological}

\begin{proposition}[共鳴状態の記述的意味]
意味的分散 $\text{Var}_E(\Delta M)$ および不確実性と価値発散指示器 $D_U, D_V$ がそれぞれそれらの閾値を下回る場合、$\text{共鳴}(t)$ が満たされる。これはセクション4の定義の直接的な言い換えである。その存在は具体的なモデル係数の割り当てにおいてのみ非自明となる。
\end{proposition}

\subsection{既存理論との相違}

Resonanceverse は単なる「非収束有界軌道」の言い直しや「意味は固定されていない」という主張ではない。その識別的特徴は差異の持続から開始し、そうした差異がいかに意味、不確実性、および価値の複数の条件下で制約され、制度的に扱う可能な構造を形成することができるかを問うことである。

固定点収束モデルでは、安定性はポイントに接近する状態として理解される。そのような枠組みの下では、差異は最終的に削減されるか消滅されるべき何かである。対照的に、本論文は完全な恒等性を前提としない。むしろそれは差異の持続それ自体を構造的条件として扱う。

古典的なコンセンサス形成理論では、中心的概念は一致、平均、またはエージェント状態間の近接である。差異はしばしばコンセンサスプロセス中に減少するべき何かとして扱われる。本論文ではコンセンサスを恒等性ではなく、$\Delta M$ 分布の有界で時間的に安定した形への修正として定義する。

有界軌道とアトラクタについての理論はその軌道の領域内への制約に対処する。しかし、それらは通常、その有界性によって保持されている意味的差異、価値的差異、または不確実性的差異が何かを問わない。また、どの差異が制度的に制約されるかを問わない。Resonanceverse の識別性はそれは有界性そのものにあるのではなく、有界性にその意味を与えるものにある。

カオスおよび複雑性理論は有界だが予測困難な振る舞いに対処する。そのような枠組みでは、差異はシステムの非線形振る舞いとして現れる。この研究が対処するのは、単なる複雑性ではなく、制約可能な $\Delta M$ 構造であり、意味、不確実性、および価値条件を同時に満たしている。

ポスト構造主義の意味理論は非固定性、\textit{différance}、および誤配に対する強力な洞察を提供する。本研究はそうした洞察を否定しない。むしろそれを $\Delta M$ ネットワークおよび共鳴条件として形式化する。従って Resonanceverse の非収束安定性は「有界非収束」ではなく、\textit{条件付き非収束}として位置付けられる。言い換えれば、この理論は単なる有界軌道ではなく、意味的差異、不確実性的差異、および価値的差異が制度的に修正される状態に対処する。


\section{哲学的考察：形式的不完全性との方法論的類推}

本セクションは数学的導出ではなく、\textit{方法論的補助線}を提示する。本論文の形式的内容はセクション2～4に含まれており、ここで述べる類推に依存しない。従って、本セクションは Resonanceverse の証明部分ではなく、非収束安定性という概念をより広い知的文脈に置くための解釈的注釈である。

形式理論のみを必要とする読者は本セクションを省略してもよい。本セクションの役割は、形式体系における「内部で完結しないが、外部からは構造的に評価されうるもの」という問題設定と、意味動力学における「収束しないが、制約下で安定しうるもの」という問題設定との間に、限定された構造的対応を示すことである。

\subsection{方法論的限定}

\textbf{方法論的限定。} ゲーデルとの対応は技術的同型ではない。Kurt Gödel の不完全性定理は、十分に強い形式体系において、体系内では証明も反証もできない命題が構成されうることを示す証明論上の結果である。一方、本論文の Resonanceverse は、意味的差異の非収束性と有界性を扱う動力学的・関係論的枠組みである。

従って、ここでの対応は「証明論的決定不可能性」と「意味的非収束性」を同一視するものではない。また、Resonanceverse がゲーデル型の不完全性定理を意味空間に対して証明する、という主張でもない。類推が有効であるのは、あくまで次の一点に限られる。

\begin{quote}
ある体系が内部で完全な閉鎖や最終的な自己同一化を達成できない場合でも、その外部または上位の観測・制約層において、整合性、持続性、構造性が評価されうる。
\end{quote}

この意味で、ゲーデル類推は証明の道具ではなく、\textit{完結性と安定性を分離して考えるための概念的装置}である。

\subsection{限定された対応関係}

ゲーデルの結果は、形式体系の内部における証明可能性と、その体系の外部から見た真理評価とを区別する契機を与えた。体系が自己の内部で全てを閉じることができないとしても、それは直ちに無秩序や無意味を意味しない。むしろ、内部的完結性の限界は、外部的評価やメタ水準の必要性を示す。

Resonanceverse においても、意味状態は単一の固定点へ収束しない。エージェント間の意味的差異
\[
\Delta M_{ij}(t)
\]
は完全には消滅せず、意味的同一性は達成されない。しかし、非収束性は不安定性と同義ではない。差異が有界領域内に保たれ、意味・不確実性・価値に関する制約を満たすとき、そこには共鳴として記述される構造的安定性が生じる。

この限定された対応は、次のように整理できる。

\begin{table}[h]
\centering
\begin{tabular}{lll}
\hline
\textbf{観点} & \textbf{形式体系} & \textbf{意味動力学} \\
\hline
内部的限界 & 体系内での決定不可能性 & 意味状態の非収束性 \\
外部的評価 & メタ水準からの真理評価 & 制約層からの共鳴評価 \\
安定性の形 & 完全性ではなく整合性 & 同一性ではなく有界な差異分布 \\
注意点 & 証明論的命題 & 動力学的・関係論的状態 \\
\hline
\end{tabular}
\end{table}

この表は厳密な対応写像ではない。むしろ、二つの領域が共有する抽象的構図を示している。すなわち、内部での完全な閉鎖が不可能であっても、上位の制約や観測を導入することで、構造的安定性を評価できるという構図である。

\subsection{真理ではなく、構造的持続性}

この類推から得られる最も重要な帰結は、真理概念の安易な拡張ではない。本論文は、形式体系における真理をそのまま意味動力学へ移植しようとするものではない。むしろ重要なのは、\textit{終端点としての真理}ではなく、\textit{非収束過程の中で持続する構造}を考えることである。

従って、本枠組みにおいて「真理」に相当するものがあるとすれば、それは単一命題や固定された意味内容ではない。それは、時間の中で変動しながらも、一定の制約条件を満たし続ける関係的パターンである。

\begin{quote}
意味的安定性は、差異の消滅ではなく、差異が評価可能な範囲に保たれることによって成立する。
\end{quote}

この定式化は、真理を固定点として扱うのではなく、意味・不確実性・価値の複数制約を満たす関係的持続性として読む可能性を示す。ただし、それは形式論理における真理概念の代替ではない。あくまで、マルチエージェント環境における意味的安定性を理解するための哲学的拡張である。

\subsection{制度設計への含意}

この方法論的類推は、制度設計に対しても示唆を持つ。制度はしばしば、意味や価値判断を一意に固定する装置として理解される。しかし Resonanceverse の観点では、制度は差異を消去するための機構ではなく、差異の分布を観測可能かつ扱いうる範囲に保つための制約構造である。

このとき、制度の目的は完全な一致を達成することではない。むしろ、異なる主体の間に残り続ける $\Delta M$ を、破局的発散や疑似的固定へ陥らせず、再解釈・再調整・再交渉が可能な範囲に保つことである。制度は意味の終着点ではなく、意味の非収束性を引き受けるための器である。

人工知能システムにおいても同様である。固定された最適解への収束だけを目指す設計は、局所的な疑似解への固定、外部擾乱への脆弱性、または価値判断の硬直化を招くことがある。非収束安定性に基づく設計は、変動を完全には抑圧せず、むしろ変動を観測・制約・更新可能な構造として扱う。これは、RDE（Resonant Deviation Evaluator）のような評価器を制度的・アーキテクチャ的に位置付ける基礎にもなる。

以上のように、ゲーデル類推は Resonanceverse の根拠ではない。それは、非完結性を欠陥としてではなく、上位の制約設計を必要とする構造条件として読むための補助線である。意味、制度、人工知能のいずれにおいても、重要なのは完全な閉鎖ではなく、閉じきれないものをどのように持続可能な構造へ編み込むかである。

\section{説明的数値シミュレーション}

理論的枠組みを例示するため、マルチエージェントシミュレーションシステムを実装し、セクション2～4および反証可能性条件に対応する四つのフェーズの実験を実施した。これらは決定的な経験的証明ではなく、明示的に指定された動力学の下で本枠組みがどのように振る舞うかを示す構成的例示である。完全な実装は \url{https://github.com/zyx-corporation/resonanceverse_sim} で利用可能である。

\subsection{Phase 1: 二エージェント基礎検証}

\textbf{目的：} スケール分離、セクション2の二エージェント例、および公理的基礎を説明的に確認する。

\textbf{セットアップ：} セクション2.4で指定されたパラメータを持つ $\mathbb{R}^2$ の二エージェントシステム。シミュレーション時刻：$T_{\max} = 100$ 時間単位、時間刻み $\Delta t = 0.01$（10,000 積分ステップ）。

\textbf{結果：}
\begin{itemize}
\item \textbf{非収束公理：} $\limsup_{t \to \infty} \|\Delta M_{12}(t)\| = 0.1009 > 0$ $\checkmark$
\item \textbf{有界性：} $\max_t \|\Delta M_{12}(t)\| = 0.1416 < \varepsilon_1 = 0.15$ $\checkmark$
\item \textbf{緩やかな変化：} $\max_t |d\|\Delta M_{12}\|/dt| = 0.0543$ は $\delta = 0.05$ をわずかに上回る。この条件は閾値に敏感であり、たとえば $\delta = 0.06$ とすれば満たされる。
\end{itemize}

\textbf{結論：} Phase 1 は、意味的差異が有界に留まる一方でゼロへ収束しないようなレジームを構成できることを示す。ただし、これは一般定理の数値的証明ではなく、モデル内での説明的例示である。補足図1を参照。

\subsection{Phase 2: マルチエージェント・コンセンサス形成}

\textbf{目的：} セクション3のコンセンサス定義および制度的制約効果を説明的に確認する。

\textbf{セットアップ：} $n = 10$ エージェント、小世界ネットワーク・トポロジー（$k = 4$、$p_{\text{rewire}} = 0.1$）。三つのシナリオ：（A）制約なし、（B）制度的制約（$\text{Var}_E < \theta_M = 0.2$）、（C）$t = 50$ での外部擾乱。

\textbf{結果：}
\begin{itemize}
\item \textbf{シナリオ A（無制約）：} コンセンサス達成 $t = 385$。
\item \textbf{シナリオ B（制約）：} コンセンサス達成 $t = 323$（42\% 高速化）。
\item \textbf{シナリオ C（擾乱）：} システムは 150 ステップ内で共鳴状態に復帰。
\end{itemize}

\textbf{結論：} 制度的制約は、この構成されたシミュレーション内で $\Delta M$ ネットワーク分布を調整しうることを示す。これはセクション3の制度設計枠組みに対する説明的支持である。補足図2を参照。

\subsection{Phase 3: 共鳴条件検証}

\textbf{目的：} 共鳴条件の同時満足可能性を説明的に確認する。

\textbf{セットアップ：} $n = 10$ エージェント、拡張状態空間 $\mathbb{R}^3$（意味 + 不確実性 + 価値）、$T_{\max} = 200$ ステップ。共鳴閾値：$\theta_M = 0.2$、$\theta_U = 0.15$、$\theta_V = 0.25$。

\textbf{閾値選定：} 閾値は無制約系（Phase 2、シナリオ A）のベースライン統計から導出された。各量 $Q \in \{\text{Var}_E, D_U, D_V\}$ について、閾値は $\theta_Q = \bar{Q}_{\text{baseline}} + 2\sigma_Q$ と設定した。ここで $\bar{Q}$ と $\sigma_Q$ は無制約実行の最終 100 タイムステップにおける $Q$ の平均と標準偏差である。これにより共鳴は非自明な条件となる：構造的協調なしには自動的に満足されない。

\textbf{手法：}

\textit{状態表現。} 各エージェント $i$ は状態ベクトル $[M_i^{(\text{sem})}, M_i^{(\text{unc})}, M_i^{(\text{val})}] \in \mathbb{R}^3$ を持ち、それぞれ意味、不確実性、価値の成分を表す。三つの成分は結合動力学の下で進化する：$M_i^{(\text{sem})}$ はセクション2で定義された相互作用関数 $f$ に従い、$M_i^{(\text{unc})}$ は入力 $\Delta M$ 信号の局所分散に比例して更新され、$M_i^{(\text{val})}$ は隣接する価値ベクトルのコサイン類似度により駆動される勾配規則に従う。

\textit{共鳴評価。} 各タイムステップ $t$ で三つの制約関数 $C_M(t), C_U(t), C_V(t)$ が定義3に従って評価される。共鳴は二値変数として記録される：$\text{共鳴}(t) = 1$ は三つの制約が全て同時に成立する場合かつその場合に限る。

\textit{リアプノフ関数。} 集約リアプノフ関数 $V(t) = \frac{1}{2}\sum_i \|M_i(t)\|^2$ が各タイムステップで計算される。その時間微分は有限差分で推定される：$\dot{V}(t) \approx (V(t+\Delta t) - V(t))/\Delta t$。下降率は $\dot{V}(t) < 0$ であるタイムステップの割合として報告される。

\textit{条件独立性。} Phase 3 中、ペアワイズ・ピアソン相関 $\text{corr}(C_M, C_U)$、$\text{corr}(C_M, C_V)$、$\text{corr}(C_U, C_V)$ が制約満足の二値時系列上で計算される。0.7 を超える相関は二つの制約がシステム状態の独立な側面を測定していないことを示す。閾値 0.7 は $R^2 = 0.49$ に対応し、一方の制約が他方の分散の半分以上を説明することを意味する。

\textbf{結果：}
\begin{itemize}
\item \textbf{共鳴達成：} タイムステップの 98.6\% が $C_M \land C_U \land C_V$ を満たす。
\item \textbf{リアプノフ下降：} $V(t) = \frac{1}{2}\sum_i \|M_i\|^2$ は 58.9\% 減少。
\item \textbf{条件独立性（動力学的）：} $C_M$ と $C_V$ 間の相関：0.609。
\item \textbf{条件独立性（統計的、Phase 4）：} 相関 0.022。
\end{itemize}

\textbf{結論：} 共鳴条件は、この設定において同時満足可能であり、安定した非収束構造に類する振る舞いを生成する。これは定義群の内的整合性を例示するが、命題群の一般的証明ではない。補足図3を参照。

\subsection{Phase 4: 反証可能性プロトコル検証}

\textbf{目的：} 反証可能性条件を診断プロトコルとして例示する。

\textbf{手法：}

\textit{プロトコル1（非収束公理）。} 反証対象は「持続的相互作用の下で $\|\Delta M_{ij}(t)\| \to 0$ となるドメインが存在する」である。これを検証するため、三つの意味的ドメイン（法的、科学的、創造的）にわたり $t_{\max} = 500$ ステップのシミュレーションを、ドメインあたり 100 回の独立試行で実行する。初期条件は $[-1, 1]^d$ から一様に抽出される。全エージェント対に対し $\|\Delta M_{ij}(t_{\max})\| < 10^{-6}$ の場合に収束と判定する。収束率（収束した試行の割合）が一次指標である。

\textit{プロトコル3（条件独立性）。} 反証対象は「$C_M, C_U, C_V$ は相互に冗長であり、単一のスカラに縮約可能である」である。これを検証するため、$\mathbb{R}^3$ で 1000 個のエージェント状態ベクトルを一様にランダム生成する。各状態に対し三つの制約述語を評価し、長さ 1000 の三つの二値ベクトルを生成する。ペアワイズ・ピアソン相関を計算する。独立性棄却閾値は $\max(\text{corr}) > 0.7$ であり、$R^2 > 0.49$ に対応する。これを超える場合、二つの制約が分散の半分以上を共有し、論理的独立性の主張が弱められる。

\textbf{プロトコル1（非収束公理）：}
\begin{itemize}
\item 複数のドメイン（法的、科学的、創造的）にわたり 100+ 試行を実施。
\item 収束率を測定：$\|\Delta M(t_{\max})\| < 10^{-6}$ （$t_{\max} = 500$ ステップ後）。
\item \textbf{結果：} 300 全試行にわたり 0\% 収束。この構成された試行群では収束は観測されなかった。これはシミュレーションモデル内では示唆的であるが、一般的な経験的証明ではない。
\end{itemize}

\textbf{プロトコル3（条件独立性）：}
\begin{itemize}
\item 1000 ランダムエージェント状態を生成。
\item ペアワイズ相関を測定：$\text{corr}(C_M, C_U)$、$\text{corr}(C_M, C_V)$、$\text{corr}(C_U, C_V)$。
\item 独立性棄却閾値：$\max(\text{corr}) > 0.7$。
\item \textbf{結果：} $\max(\text{corr}) = 0.022 \ll 0.7$。このサンプリングプロトコルの下では低い相関が観測された。これは選択された符号化における非冗長性を支持する。
\end{itemize}

\textbf{結論：} 反証可能性プロトコルは、Resonanceverse を検査可能な理論枠組みとして扱うための診断手順を例示する。非収束公理はこのシミュレーション内では支持的に振る舞い、共鳴条件の非冗長性も示唆されるが、いずれも普遍的な経験的証明ではない。補足図4を参照。

\section{結論}

本研究は Resonanceverse を提示した。これは意味を静的な表現ではなく、エージェント間の相互作用を通じて生成される動的構造として再定義する理論的枠組みである。本枠組みでは、意味は固定的なオブジェクトではなく、時間の中で継続的に変動する状態量であり、その変換はエージェント間の差異、すなわち $\Delta M$ によって記述される。

セクション2は意味状態空間、意味的差異 $\Delta M_{ij}(t) := M_j(t) - M_i(t)$、および関係を含む公理系を導入した。それは $\Delta M$ が誤配そのものではなく、誤配プロセスから生じる状態的差異であることを明確にした。本セクションはまた非収束公理を実用的 $\varepsilon$ スケール・コンセンサスおよび共鳴からスケール分離を通じて区別し、厳密な恒等が原理的には達成されないことを明確にした。

セクション3 はフレームワークをマルチエージェント・システムに拡張し、意味論的構造を意味的差異のネットワークとして記述した。それはコンセンサスを恒等性ではなく、意味的差異の分布が有界で時間的に安定する状態として再定義した。従ってコンセンサスは差異の消滅ではなく、差異の構造的修正として示された。

セクション4 は非収束安定性を非収束状態の安定性概念として導入し、共鳴を $C_M, C_U, C_V$ の同時満足として定義した。それは意味が収束することなく安定した構造を形成することがあり得ることを明らかにする理論的足場を整理した。振幅有界性と非収束性が仮定1などの条件下で共存できることを示した。

セクション5 は、形式的不完全性との慎重に限定された方法論的類推を通じて、フレームワークをより広い理論的文脈に位置付けた。この類推は数学的根拠としてではなく、完結性と安定性を分離して考えるための概念的補助線として用いられた。そこでは、終端点としての真理ではなく構造的持続性が強調され、非収束安定性は RDE のような評価器を含む制度的・アーキテクチャ的制約設計へ接続された。

セクション6 は四つのフェーズにわたる説明的数値シミュレーションを提示した。実験は、明示的に指定された動力学の下で、有界な非収束的振る舞いを構成できること、制度的制約効果を観測できること、$C_M \land C_U \land C_V$ に類する共鳴条件満足を得られること、および反証可能性を志向した診断を定式化できることを示した。これらの結果は提案されたモデル化枠組みの整合性を支持するが、非収束公理の普遍的な経験的証明ではない。

\subsection{意味の再定義}

これらの結果から、本論文は意味の以下の再定義を提供する。

\textit{意味は主体に内在する固定的内容ではない。それはエージェント間の差異のネットワークとして存在し、そのような差異の時間的累積を通じて形成される。その安定性は収束によってではなく、差異を保持しながら制約された範囲内で持続することによって成立する。}

この再定義はセマンティクスだけでなく、コンセンサス形成、制度設計、および人工知能にも含意を持つ。特に、収束志向モデルを非収束安定性で置き換えることは、外部擾乱に対する堅牢性と誤った固定の回避に対して有用な枠組みを提供することができる。

\subsection{反証可能性と限定条件}

本研究は意味が原理的に完全な恒等性に収束しないという非収束公理に根ざしている。しかし、この公理は全ての対象に無条件に適用される実証的事実として主張されない。従って、この理論は以下の条件下で改訂または棄却されることがある。

\textbf{条件1：} 複数の意味的ドメイン全体にわたって、符号化された意味的差異 $\Delta M_{ij}(t)$ が十分に長い相互作用の後、安定的にゼロに収束することが反復可能に示される場合、非収束公理を一般原理として維持することはできず、そのスコープを限定する必要がある。

\textbf{条件2：} $\Delta M$ の分布が、制度的制約または文脈的親和性係数 $\alpha_{ij}(t)$ なしに、一般的に有界で安定することが示される場合、$\Delta M$ 分布への制約として制度を位置付ける必要性は弱まる。

\textbf{条件3：} 意味、不確実性、および価値の三つの条件——$C_M, C_U, C_V$——が独立として分離されず、代わりに単一のスカラ指標に削減されることができる場合、その論理的論理積として共鳴を定義する有意性は減少する。

\textbf{条件4：} $\Delta M$ が観測可能または推定可能な量として安定的に構成されない場合、理論は形式モデルとしての堅牢性を失う。

\textbf{条件5：} 共鳴条件を満たす状態が、実際のマルチエージェント・システムで安定性ではなく抑圧、均質化、または病的閉鎖を生成する場合、共鳴は望ましい安定性ではなく疑似安定性として再定義されなければならない。

これらの限定条件は理論の単なる弱点ではなく、将来の研究課題でもある。Resonanceverse は完成した閉じられたシステムではない。それは $\Delta M$ の観測、制度的制約の設計、および共鳴条件の定量化を通じて検査されなければならない開放的な理論的枠組みである。

\subsection{将来の仕事}

将来の仕事はここで提示された枠組みの数学的精密化を含む。具体的には、$\Delta M$ を確率過程として形式化し、共鳴条件を定量的に評価し、非収束安定性の存在と一意性に関する定理を導出することが必要であろう。追加的なマルチエージェント・シミュレーションおよび実際の意味的ドメインにおける実験的検証を通じて理論をテストすることもまた重要である。

制度および人工知能システムへの理論応用では、中心的設計問題は $\Delta M$ の分布をいかに制御するか、およびどの条件の下で共鳴が維持されるかであろう。その点で本研究は意味動力学に根ざした新たな設計原則を提供する。

本論文は意味を収束的オブジェクトではなく、非収束的でかつ構造的に安定したプロセスとして理解するための理論的基礎を提示した。この視点は意味の再構築だけでなく、複雑なマルチエージェント・システムにおける安定性の理解へも新しい枠組みを提供する。

簡潔に、本研究の立場は以下のように述べられることがある。

\textit{意味は恒等性を通じて生じない。それは差異を保持しながら持続する構造として現れる。この原則は セマンティクスの再構築のためだけでなく、複雑なマルチエージェント・システムにおける安定性の理解のためにも指針を提供する。}

\begin{thebibliography}{99}

\bibitem{azuma1998} H. Azuma, \textit{Ontological, Postal: On Jacques Derrida} [Sonzaironteki, Yubinteki: Jacques Derrida ni tsuite], Shinchosha, 1998.

\bibitem{berger1966} P. L. Berger and T. Luckmann, \textit{The Social Construction of Reality: A Treatise in the Sociology of Knowledge}, Doubleday, 1966.

\bibitem{degroot1974} M. H. DeGroot, ``Reaching a consensus,'' \textit{Journal of the American Statistical Association}, vol. 69, no. 345, pp. 118--121, 1974.

\bibitem{derrida1967} J. Derrida, \textit{De la Grammatologie}, Les Éditions de Minuit, 1967.

\bibitem{derrida1972} J. Derrida, ``La différance,'' in \textit{Marges de la Philosophie}, pp. 1--29, Les Éditions de Minuit, 1972.

\bibitem{godel1931} K. Gödel, ``Über formal unentscheidbare Sätze der Principia Mathematica und verwandter Systeme I,'' \textit{Monatshefte für Mathematik und Physik}, vol. 38, pp. 173--198, 1931.

\bibitem{khalil2002} H. K. Khalil, \textit{Nonlinear Systems} (3rd ed.), Prentice Hall, 2002.

\bibitem{landauer1997} T. K. Landauer and S. T. Dumais, ``A solution to Plato's problem: The latent semantic analysis theory of acquisition, induction, and representation of knowledge,'' \textit{Psychological Review}, vol. 104, no. 2, pp. 211--240, 1997.

\bibitem{mikolov2013} T. Mikolov, K. Chen, G. Corrado, and J. Dean, ``Efficient estimation of word representations in vector space,'' arXiv:1301.3781, 2013.

\bibitem{moreau2005} L. Moreau, ``Stability of multiagent systems with time-dependent communication links,'' \textit{IEEE Transactions on Automatic Control}, vol. 50, no. 2, pp. 169--182, 2005.

\bibitem{north1990} D. C. North, \textit{Institutions, Institutional Change and Economic Performance}, Cambridge University Press, 1990.

\bibitem{olfati2007} R. Olfati-Saber, J. A. Fax, and R. M. Murray, ``Consensus and cooperation in networked multi-agent systems,'' \textit{Proceedings of the IEEE}, vol. 95, no. 1, pp. 215--233, 2007.

\bibitem{shannon1948} C. E. Shannon, ``A mathematical theory of communication,'' \textit{The Bell System Technical Journal}, vol. 27, no. 3, pp. 379--423; vol. 27, no. 4, pp. 623--656, 1948.

\bibitem{sontag1989} E. D. Sontag, ``Smooth stabilization implies coprime factorization,'' \textit{IEEE Transactions on Automatic Control}, vol. 34, no. 4, pp. 435--443, 1989.

\bibitem{strogatz2015} S. H. Strogatz, \textit{Nonlinear Dynamics and Chaos: With Applications to Physics, Biology, Chemistry, and Engineering} (2nd ed.), Westview Press, 2015.

\end{thebibliography}

\end{document}
